\appendix
\chapter{Annexes}

\section{Extraits de code sélectionnés}

L'extrait ci-dessous illustre, sous une forme simplifiée, la logique de fusion des scores de l'ensemble de détection d'anomalies (module P4), avec renormalisation des poids lorsqu'un détecteur est indisponible :

\begin{verbatim}
def fuse_anomaly_scores(scores: dict, weights: dict) -> float:
    available = {k: v for k, v in scores.items() if v is not None}
    total_weight = sum(weights[k] for k in available)
    if total_weight == 0:
        return 0.0
    return sum(scores[k] * weights[k] for k in available) / total_weight
\end{verbatim}

L'extrait suivant illustre le principe du calcul du score de disponibilité opérationnelle (module P7), combinant plusieurs signaux pondérés :

\begin{verbatim}
def compute_readiness_score(health, inventory, shortage, recency):
    return (
        0.40 * health
        + 0.30 * inventory
        - 0.20 * shortage
        + 0.10 * recency
    )
\end{verbatim}

\section{Diagrammes détaillés}

Le diagramme de classes complet ainsi que le diagramme de séquence du pont ML-RAG (invocation parallèle du service RAG et du microservice ML, fusion du contexte avant génération de la réponse) sont disponibles dans les spécifications techniques du projet.

\section{Tableaux de métriques complets}

\begin{table}[h]
\centering
\begin{tabularx}{\textwidth}{|l|X|}
\toprule
\textbf{Modèle} & \textbf{Description} \\
\midrule
P1 & Probabilité de panne (classification binaire) \\
\midrule
P2 & Classification du type de panne \\
\midrule
P3 & Estimation de la durée de vie utile restante (MAE = 14,95 cycles, indice de concordance = 0,620) \\
\midrule
P4 & Détection d'anomalie comportementale (ROC-AUC de référence = 0,83) \\
\midrule
P5 & Priorisation des ordres de travail \\
\midrule
P6 & Planification de la maintenance \\
\midrule
P7 & Coordination des pièces détachées (prévision de la demande, alertes, approvisionnement guidé) \\
\bottomrule
\end{tabularx}
\caption{Récapitulatif des modèles du pipeline d'apprentissage automatique}
\label{tab:metriques-completes-fr}
\end{table}
